\chapter{Планирование эксперимента}

\section{Постановка задачи и исследовательские гипотезы}

Настоящая глава посвящена планированию экспериментального исследования. Основной научный вклад работы сосредоточен на уровне \emph{групповой агрегации}: предлагается семейство audio-aware attention-агрегаторов и строится экспериментальный протокол его сравнения с современными бейзлайнами групповых рекомендаций (AGREE, GroupIM) и классическими стратегиями агрегации.

В целях методологической чистоты экспериментальный дизайн работы намеренно сужен по сравнению с полной четырёхэтапной воронкой главы~1. Полная воронка является архитектурным ориентиром продуктовой части диплома; в технической части она используется как фон, а в качестве объекта эксперимента выделяется единственный компонент~--- групповой агрегатор. Конкретно это означает: (а)~на этапе персонального последовательного ранжирования зафиксирована одна архитектура~--- gSASRec, без ablation против альтернатив; (б)~финальный CatBoost-этап не включается в групповой режим, чтобы избежать сдвига распределений на его входе; (в)~групповой агрегатор работает непосредственно поверх скоров последовательной модели, отображая объединение персональных топ-кандидатов в финальный групповой список. Такое сужение увеличивает статистическую мощность сравнений и обеспечивает однозначную интерпретацию различий между моделями.

\subsection{Формализация задачи персональных рекомендаций}

Пусть $\mathcal{U} = \{u_1, \ldots, u_{|\mathcal{U}|}\}$~--- множество пользователей, $\mathcal{I} = \{i_1, \ldots, i_{|\mathcal{I}|}\}$~--- каталог музыкальных треков. История взаимодействий пользователя $u \in \mathcal{U}$ представляется упорядоченной во времени последовательностью
\begin{equation}
    \mathcal{S}_u = \left( (i_1^{(u)}, t_1^{(u)}), (i_2^{(u)}, t_2^{(u)}), \ldots, (i_{n_u}^{(u)}, t_{n_u}^{(u)}) \right),
\end{equation}
где $i_k^{(u)} \in \mathcal{I}$~--- трек, прослушанный пользователем $u$ на шаге $k$, а $t_k^{(u)}$~--- момент времени, причём $t_1^{(u)} < \ldots < t_{n_u}^{(u)}$.

Персональная модель $f_\theta$ с параметрами $\theta$ по истории $\mathcal{S}_u^{(<t)} = \left( i_1^{(u)}, \ldots, i_{t-1}^{(u)} \right)$ оценивает релевантность трека $i \in \mathcal{I}$ на следующем шаге:
\begin{equation}
    \hat{s}_{u, i}^{(t)} = f_\theta\!\left( \mathcal{S}_u^{(<t)}, i \right).
\end{equation}

Практическая реализация $\hat{s}$ в эксперименте строится как композиция трёх этапов: генерация кандидатов через ALS с индексом приближённого ближайшего соседа ($|\mathcal{I}| \to 10^4$), лёгкое pre-ranking лёгким CatBoost ($10^4 \to 10^3$) и последовательное ранжирование архитектурой gSASRec~\cite{gsasrec2023} ($10^3 \to 10^2$). Четвёртый, финальный этап тяжёлого CatBoost-ранжирования, описанный в архитектурной части главы~1, в техническом эксперименте настоящей работы \emph{не используется}: он создаёт сдвиг распределений между обучением (на персональных входах) и инференсом (на групповых входах), что вносит дополнительный неконтролируемый источник вариации. Применение этапа~4 относится к продуктовой части диплома.

Выбор gSASRec в качестве модели последовательного ранжирования зафиксирован и не подвергается ablation. Обоснование: согласно литературе~\cite{gsasrec2023}, gSASRec систематически превосходит классические трансформерные бейзлайны (SASRec, BERT4Rec) на больших каталогах за счёт функции потерь gBCE, компенсирующей смещение от отрицательного семплирования; YAMBDA с $|\mathcal{I}| \approx 9{,}4 \cdot 10^6$ относится к классу больших каталогов. Сравнение с альтернативными архитектурами (в частности, Mamba4Rec~\cite{liu2024mamba4rec}) представляет самостоятельный интерес, но выходит за рамки задачи о групповой агрегации и в случае необходимости выносится в приложение работы.

\subsection{Формализация задачи групповых рекомендаций и предлагаемая модификация агрегатора}

Пусть $G = \{u_1, \ldots, u_{|G|}\} \subseteq \mathcal{U}$~--- эфемерная группа пользователей, сформированная на время совместного потребления музыки (посетители заведения, участники мероприятия). Задача групповых рекомендаций состоит в построении списка $\mathrm{Rec}_K(G, t)$ длины $K$, который в совокупности удовлетворяет всех членов группы.

Используется двухуровневый подход. На первом уровне для каждого пользователя $u \in G$ вычисляются персональные оценки $\hat{s}_{u, i}^{(t)}$ с помощью зафиксированного per-user скорера. На втором уровне применяется функция агрегации $\mathcal{A}$, отображающая множество персональных оценок в групповую:
\begin{equation}
    \hat{s}_{G, i}^{(t)} = \mathcal{A}\!\left( \left\{ \hat{s}_{u, i}^{(t)} \right\}_{u \in G}, \, \mathbf{a}_i, \, \{\mathbf{a}_j\}_{j \in \mathcal{S}_u} \right),
    \qquad
    \mathrm{Rec}_K(G, t) = \underset{i \in \mathcal{I}}{\mathrm{argtop}\text{-}K}\, \hat{s}_{G, i}^{(t)},
\end{equation}
где $\mathbf{a}_i \in \mathbb{R}^{d_a}$~--- аудиоэмбеддинг трека $i$, предоставляемый датасетом YAMBDA~\cite{yambda2025}.

В качестве базовых стратегий агрегации рассматриваются классические: среднее (Average), стратегия наименьшего недовольства (Least Misery) и максимального удовольствия (Max Pleasure)~\cite{cao2018agree}. В качестве обучаемого бейзлайна~--- attention-агрегатор AGREE~\cite{cao2018agree}, использующий обучаемые ID-эмбеддинги треков $\mathbf{e}_i$.

\paragraph{Предлагаемая модификация: audio-aware attention-агрегатор.}
Ключевой недостаток AGREE в музыкальном домене~--- зависимость от обучаемого ID-эмбеддинга $\mathbf{e}_i$, качество которого резко падает на редких и новых треках. Мы предлагаем заменить $\mathbf{e}_i$ на аудиоэмбеддинг $\mathbf{a}_i$, доступный для любого трека из каталога независимо от частоты его прослушиваний. Это позволяет агрегатору работать на cold-start треках без дообучения и вводит в модель сигнал, отсутствующий в существующих работах по групповым рекомендациям.

В работе исследуются три варианта модификации по возрастанию сложности:
\begin{enumerate}
    \item \textbf{Audio taste profile} (необучаемый бейзлайн). Для каждого пользователя строится аудиопрофиль $\bar{\mathbf{a}}_u = \frac{1}{|\mathcal{S}_u|} \sum_{j \in \mathcal{S}_u} \mathbf{a}_j$. Вес пользователя определяется как softmax косинусных сходств профиля с аудиоэмбеддингом кандидата:
    \begin{equation}
        \alpha_{u, i}^{(G)} = \frac{\exp\!\left( \tau \cdot \cos(\bar{\mathbf{a}}_u, \mathbf{a}_i) \right)}{\sum_{u' \in G} \exp\!\left( \tau \cdot \cos(\bar{\mathbf{a}}_{u'}, \mathbf{a}_i) \right)},
    \end{equation}
    где $\tau$~--- температурный параметр. Обучаемых параметров агрегатора почти нет~--- это базовая точка сравнения для оценки вклада архитектуры.

    \item \textbf{Audio-AGREE} (независимый attention). Веса пользователей вычисляются обучаемой скоринговой функцией $\phi_\psi$ с параметрами $\psi$ поверх аудиопрофиля и аудиоэмбеддинга кандидата:
    \begin{equation}
        \alpha_{u, i}^{(G)} = \frac{\exp\!\left( \phi_\psi(\bar{\mathbf{a}}_u, \mathbf{a}_i) \right)}{\sum_{u' \in G} \exp\!\left( \phi_\psi(\bar{\mathbf{a}}_{u'}, \mathbf{a}_i) \right)}.
    \end{equation}
    Архитектура повторяет AGREE, но оперирует аудиоэмбеддингами вместо обучаемых ID-эмбеддингов.

    \item \textbf{Group cross-attention} (взаимодействие членов группы). Кандидат $\mathbf{a}_i$ используется как запрос (query), аудиопрофили членов группы~--- как ключи и значения (keys, values) в многоголовом cross-attention блоке. Это позволяет весу одного пользователя учитывать наличие других пользователей в группе, а не только их собственные профили.
\end{enumerate}
Во всех трёх вариантах итоговая групповая оценка вычисляется как взвешенное среднее персональных скоров:
\begin{equation}
    \hat{s}_{G, i}^{(t)} = \sum_{u \in G} \alpha_{u, i}^{(G)} \hat{s}_{u, i}^{(t)}.
\end{equation}

Сравнение этих трёх вариантов между собой и с классическими бейзлайнами (AVG, Least Misery, Max Pleasure, ID-based AGREE) составляет ядро экспериментального исследования работы.

\paragraph{Замечание об анонимизации данных.}
Датасет YAMBDA полностью анонимизирован: идентификаторы треков являются внутренними кодами Яндекса, а аудиоэмбеддинги предоставляются без какой-либо семантической разметки (жанр, исполнитель, настроение). Это исключает возможность семантической верификации эмбеддингов. В качестве замены в настоящей работе применяется \emph{поведенческая валидация} как методологическая sanity-проверка: до основного эксперимента подтверждается, что геометрическая близость в пространстве аудиоэмбеддингов коррелирует с поведенческими паттернами в данных (ко-встречаемость треков в сессиях, сходство временных профилей прослушивания). Детали процедуры приведены в разделе~2.2.5; её положительный результат рассматривается как предусловие осмысленности гипотез H1--H2, но не как самостоятельная научная гипотеза.

Группы пользователей в YAMBDA не заданы~--- их синтез описан в разделе~2.3. Основная линия экспериментов работает с фиксированными на сеанс группами; динамический сценарий (приход и уход участников в течение сеанса) выносится в отдельный раздел, связывающий техническую часть работы с последующей продуктовой.

\subsection{Исследовательские гипотезы}

Эксперимент сосредоточен на проверке двух центральных гипотез о предложенном семействе audio-aware агрегаторов. Дополнительно ставится вспомогательный вопрос об устойчивости эффекта на различных типах групп и редких треках; методологическая корректность использования аудиоэмбеддингов в условиях анонимизации проверяется отдельной sanity-проверкой (раздел~2.2.5).

\textbf{Гипотеза H1 (преимущество audio-aware агрегации).} Предложенное семейство audio-aware attention-агрегаторов превосходит классические стратегии (Average, Least Misery, Max Pleasure), а также современные обучаемые бейзлайны для эфемерных групп~--- ID-based AGREE~\cite{cao2018agree} и GroupIM~\cite{sankar2020groupim}~--- по основным групповым метрикам настоящей работы (Group-NDCG@K, Jain@K и их сводной мере DFH@K, см. раздел~2.4). Ожидаемое преимущество обусловлено тем, что аудиоэмбеддинги дают агрегатору информацию о содержательном сходстве треков, недоступную моделям на ID-эмбеддингах и особенно ценную на редких треках.

\textbf{Гипотеза H2 (роль архитектурной сложности агрегатора).} В пределах семейства audio-aware агрегаторов качество монотонно растёт с усложнением архитектуры: \emph{group cross-attention} $>$ \emph{Audio-AGREE} $>$ \emph{audio taste profile}. Проверка данной гипотезы отвечает на вопрос, даёт ли выигрыш сам факт использования аудио или требуется именно моделирование взаимодействия членов группы.

\textbf{Дополнительный анализ (тип группы и cold-start).} В отдельных срезах оценивается, усиливается ли преимущество audio-aware агрегатора (а)~на гетерогенных группах с сильно различающимися вкусами и (б)~на множествах кандидатов с высокой долей редких треков. Этот анализ не оформляется как самостоятельная гипотеза, поскольку зависит от исхода H1, и приводится для содержательной интерпретации источников выигрыша.

Поведенческая валидация аудиоэмбеддингов (со-встречаемость треков в сессиях, поведенческая кластеризация, согласованность пользовательских профилей) рассматривается как методологическая sanity-проверка перед основным экспериментом, а не как самостоятельная научная гипотеза: её положительный результат~--- предусловие осмысленности H1--H2 в условиях анонимизированного датасета. Детали процедуры приведены в разделе~2.2.5.

\section{Датасет YAMBDA и его подготовка}

\subsection{Характеристики и состав датасета}

В качестве источника данных используется YAMBDA~\cite{yambda2025}~--- открытый датасет взаимодействий пользователей с музыкальным каталогом, опубликованный командой Яндекс Музыки в 2025~году. Датасет содержит около $4{,}79 \cdot 10^9$ событий взаимодействий, $|\mathcal{U}| \approx 10^6$ уникальных пользователей и $|\mathcal{I}| \approx 9{,}4 \cdot 10^6$ уникальных треков, что делает его одним из крупнейших публично доступных ресурсов для рекомендательных задач в музыкальном домене. Каждое событие представлено кортежем $(u, i, t, \mathrm{type}, \mathrm{is\_organic})$, где $u$~--- идентификатор пользователя, $i$~--- идентификатор трека, $t$~--- временная метка, $\mathrm{type}$~--- тип взаимодействия (прослушивание, лайк, дизлайк, пропуск), $\mathrm{is\_organic} \in \{0, 1\}$~--- флаг, отделяющий органические действия пользователя от событий, инициированных рекомендательной системой платформы.

Для каждого трека $i \in \mathcal{I}$ в датасете предоставлен предобученный аудиоэмбеддинг $\mathbf{a}_i \in \mathbb{R}^{d_a}$, полученный внутренней моделью Яндекса. Семантическая интерпретация осей эмбеддинга, а также метаданные треков (исполнитель, жанр, год выпуска, текстовое описание) в открытой версии датасета отсутствуют: идентификаторы $i$ являются анонимизированными внутренними кодами и не сопоставимы с публичными музыкальными базами. Это обстоятельство принципиально определяет процедуру валидации эмбеддингов, описанную в разделе~2.2.5.

Принципиально важно отделять YAMBDA (2025) от датасета \emph{Yandex Music Event} (2019), также публиковавшегося Яндексом и часто используемого в литературе по последовательным рекомендациям~\cite{abbattista2024sequential}: последний содержит около $4{,}77 \cdot 10^7$ событий, не включает аудиоэмбеддинги и флаг \texttt{is\_organic} и не пригоден для задач, рассматриваемых в настоящей работе.

\subsection{Фильтрация событий}

Перед построением экспериментальных выборок к исходному датасету применяется последовательность шагов фильтрации, направленных на исключение шума и сокращение объёма данных до уровня, совместимого с доступными вычислительными ресурсами.

\textbf{Тип взаимодействий.} В качестве позитивных взаимодействий рассматриваются только события прослушивания трека на сумму свыше заданного порога $\theta_l$ (доля от длительности трека). Пропуски, дизлайки и взаимодействия с пользовательским интерфейсом исключаются: они либо несут отрицательный сигнал, либо не отражают содержательного интереса к треку.

\textbf{Органические события.} Из обучающих и тестовых выборок удаляются события с $\mathrm{is\_organic} = 0$. Это решение обосновано следующим: события, инициированные собственной рекомендательной системой Яндекс Музыки, отражают не предпочтения пользователя, а поведение существующего рекомендера, и их использование привело бы к смещению оценок в пользу моделей, имитирующих поведение исходного рекомендера.

\textbf{Порог активности пользователей и треков.} Рассматриваются только пользователи с $|\mathcal{S}_u| \geq n_u^{\min}$ органическими прослушиваниями и треки с не менее чем $n_i^{\min}$ органическими прослушиваниями за период наблюдения. Конкретные значения $n_u^{\min}$ и $n_i^{\min}$ выбираются по результатам разведочного анализа таким образом, чтобы (а)~сохранить достаточную статистическую мощность экспериментов, (б)~оставить нетривиальную долю треков с малой частотой для корректной оценки cold-start сценария (раздел~2.4) и (в)~обеспечить совместимость с вычислительными ограничениями (раздел~2.7).

\textbf{Повторные прослушивания.} В отличие от ряда работ, агрегирующих повторные прослушивания одного и того же трека пользователем в одно взаимодействие, в настоящей работе повторные прослушивания сохраняются. Это обусловлено спецификой музыкального домена: повторное потребление является существенной составляющей пользовательского поведения~\cite{abbattista2024sequential, greer2025beyond} и должно быть доступно моделям как сигнал.

\subsection{Временной сплит}

В качестве протокола разделения данных на обучающую, валидационную и тестовую выборки используется \emph{глобальный временной сплит} (global timeline split). Пусть $T$~--- общая временная протяжённость датасета. Выбираются две точки во времени $t_{\mathrm{train}} < t_{\mathrm{val}} < T$, и для каждого пользователя $u$:
\begin{itemize}
    \item обучающая выборка содержит события $\{ (i, t) \in \mathcal{S}_u : t \leq t_{\mathrm{train}} \}$;
    \item валидационная~--- события $\{ (i, t) \in \mathcal{S}_u : t_{\mathrm{train}} < t \leq t_{\mathrm{val}} \}$;
    \item тестовая~--- события $\{ (i, t) \in \mathcal{S}_u : t > t_{\mathrm{val}} \}$.
\end{itemize}

Данный протокол предпочтительнее распространённой схемы leave-one-out по следующим причинам. Во-первых, он не нарушает каузальность во времени: при обучении модель не имеет доступа ни к одному событию из будущего ни одного пользователя, что воспроизводит реальную production-постановку~\cite{volokitin2023yandex}. Во-вторых, он корректно учитывает популярные тренды (выход новых треков, сезонные эффекты), оценивая способность модели обобщать на действительно новые элементы каталога, а не на случайно скрытые из истории. В-третьих, исследования воспроизводимости~\cite{petrov2022bert4rec_replicability} показывают, что leave-one-out склонен переоценивать качество последовательных моделей за счёт утечек.

Длины окон $t_{\mathrm{train}}$ и $t_{\mathrm{val}}$ подбираются так, чтобы валидационный и тестовый периоды содержали достаточно событий для устойчивой оценки метрик и при этом основной объём данных оставался в обучении.

\subsection{Подвыборка под вычислительные ограничения}

Эксперименты выполняются в среде Google Colab с доступом к графическим ускорителям A100 на ограниченное по продолжительности время. Это исключает возможность обучения моделей на полном объёме YAMBDA ($4{,}79 \cdot 10^9$ событий). Для приведения данных к допустимому размеру применяется стратифицированная подвыборка пользователей.

Пусть $\mathcal{U}_{\mathrm{filt}}$~--- множество пользователей, прошедших фильтрацию по разделу~2.2.2. Из $\mathcal{U}_{\mathrm{filt}}$ случайным образом отбирается подмножество $\mathcal{U}_{\mathrm{exp}}$ заранее выбранного размера $|\mathcal{U}_{\mathrm{exp}}| = N_u$. Стратификация проводится по двум атрибутам: (а)~общая активность пользователя, разбитая на квантили распределения $|\mathcal{S}_u|$; (б)~разнообразие потребления, оценённое как энтропия распределения времени прослушивания по поведенческим кластерам треков (раздел~2.2.5). Стратифицированный отбор позволяет сохранить в подвыборке баланс между активными и малоактивными пользователями, а также между пользователями с узкими и широкими вкусами, что критично для корректной оценки гипотезы H3.

Каталог треков ограничивается треками, встречающимися хотя бы в одной истории $\mathcal{S}_u$, $u \in \mathcal{U}_{\mathrm{exp}}$. Все случайные выборы фиксируются единым набором сидов и сохраняются в виде идентификаторов пользователей и треков, что обеспечивает побитовую воспроизводимость экспериментов.

\subsection{Поведенческая валидация аудиоэмбеддингов}

Как было отмечено в разделе~2.1, анонимизация YAMBDA исключает семантическую интерпретацию аудиоэмбеддингов. Тем не менее, для использования $\mathbf{a}_i$ в качестве сигнала в групповом агрегаторе необходимо убедиться, что геометрия эмбеддинг-пространства согласована с наблюдаемыми пользовательскими паттернами~--- то есть проверить гипотезу H4. В настоящей работе предлагается следующая процедура валидации.

\textbf{Со-встречаемость в сессиях.} Сессией считается максимальная по включению последовательность событий пользователя, для которой интервал между соседними событиями не превышает $\Delta t_s$. Для каждой пары треков $(i, j)$, встретившихся в одной сессии хотя бы одного пользователя, вычисляется их косинусное расстояние в пространстве эмбеддингов $d(\mathbf{a}_i, \mathbf{a}_j)$. Контрольная выборка формируется случайным семплированием пар треков с учётом распределения частот. Гипотеза H4 считается подтверждённой, если средняя близость со-встречающихся пар статистически значимо выше, чем у контрольной выборки (одностороннее $t$-сравнение средних с поправкой Бонферрони на множественность).

\textbf{Поведенческая кластеризация.} На случайной подвыборке треков выполняется кластеризация эмбеддингов методом $k$-means с числом кластеров, выбираемым по индексу силуэта. Каждому кластеру сопоставляется поведенческий профиль: средняя популярность входящих в него треков, доля \texttt{is\_organic}-событий, типичное время суток прослушивания, средняя длина сессии, доля повторных прослушиваний. Различия профилей между кластерами проверяются попарно методом Краскела--Уоллиса. Содержательная различимость поведенческих профилей кластеров служит дополнительным свидетельством наличия структуры в эмбеддинг-пространстве.

\textbf{Согласованность пользовательских профилей.} Для каждого пользователя $u \in \mathcal{U}_{\mathrm{exp}}$ строится аудиопрофиль $\bar{\mathbf{a}}_u$, после чего проверяется, что косинусное сходство профилей двух случайных пользователей с пересекающимися историями выше, чем у пользователей с непересекающимися историями. Это базовая sanity-проверка корректности конструкции $\bar{\mathbf{a}}_u$, используемой в audio-aware агрегаторе.

Отрицательный результат любой из перечисленных проверок означает, что геометрия аудиоэмбеддингов не несёт сигнала, релевантного задаче, и требует пересмотра роли $\mathbf{a}_i$ в архитектуре агрегатора. Подтверждение H4 обосновывает методологическую допустимость предлагаемого подхода в условиях анонимизации.

\section{Синтез групп пользователей}

\subsection{Обоснование процедуры синтеза}

Датасет YAMBDA не содержит сведений о реально существующих группах пользователей: каждое событие связано с одиночным идентификатором $u$, а информация о совместном прослушивании музыки в физически локализованных группах (заведения, мероприятия) принципиально недоступна. Это типичная ситуация в литературе по групповым рекомендациям~\cite{cao2018agree, sankar2020groupim, xu2024aligngroup}: эфемерные группы, как правило, отсутствуют в открытых датасетах и формируются искусственно из индивидуальных профилей.

Синтез эфемерных групп преследует две цели. Во-первых, обеспечить контролируемое варьирование характеристик группы (размер, степень сходства членов), что необходимо для проверки гипотезы H3 о различном поведении агрегаторов на гомогенных и гетерогенных группах. Во-вторых, исключить смещение, связанное со случайной структурой групп: при чисто случайном объединении пользователей преобладают группы со средней степенью сходства, что не позволяет различить вклад агрегатора в крайних режимах. Применяемая в работе процедура строит сбалансированный набор групп по сценариям сходства и размеру.

Формально, набор синтезированных групп определяется как
\begin{equation}
    \mathcal{G} = \left\{ G \subseteq \mathcal{U}_{\mathrm{exp}} : |G| \in \{2, 4, 8\}, \, \mathrm{type}(G) \in \{\mathrm{hom}, \mathrm{het}, \mathrm{rand}\} \right\},
\end{equation}
где $\mathrm{type}(G)$~--- сценарий формирования группы (гомогенная, гетерогенная, случайная), а $|G|$~--- размер группы. Для каждой комбинации $(|G|, \mathrm{type}(G))$ генерируется фиксированное число групп $N_G$ с разными случайными сидами, что обеспечивает статистическую устойчивость оценок.

\subsection{Мера сходства пользователей}

Сценарии формирования групп опираются на меру сходства между парами пользователей. Чтобы избежать самореференции (использования аудиоэмбеддингов одновременно как сигнала агрегатора и как критерия формирования групп), сходство пользователей определяется на основе \emph{пересечения историй прослушивания}, а не на основе $\bar{\mathbf{a}}_u$. Используется коэффициент Жаккара по множествам прослушанных треков:
\begin{equation}
    \mathrm{sim}(u, v) = \frac{|\mathcal{S}_u^{\mathrm{set}} \cap \mathcal{S}_v^{\mathrm{set}}|}{|\mathcal{S}_u^{\mathrm{set}} \cup \mathcal{S}_v^{\mathrm{set}}|},
\end{equation}
где $\mathcal{S}_u^{\mathrm{set}} \subseteq \mathcal{I}$~--- множество уникальных треков, прослушанных пользователем $u$ в обучающем периоде. Использование коэффициента Жаккара мотивировано двумя соображениями: (а)~независимость от длины истории, что важно при широком диапазоне $|\mathcal{S}_u|$ в YAMBDA; (б)~независимость от пространства эмбеддингов, что устраняет циркулярность при оценке audio-aware агрегатора.

В дополнение к $\mathrm{sim}(u, v)$ для аналитических целей вычисляется аудио-сходство $\mathrm{sim}_a(u, v) = \cos(\bar{\mathbf{a}}_u, \bar{\mathbf{a}}_v)$, не используемое при формировании групп, но применяемое в анализе результатов.

\subsection{Стратегии формирования групп}

Рассматриваются три сценария формирования групп.

\textbf{Случайные группы (random).} Группа $G$ формируется случайной выборкой $|G|$ пользователей из $\mathcal{U}_{\mathrm{exp}}$ без ограничений на $\mathrm{sim}(u, v)$. Этот сценарий служит контрольной точкой и воспроизводит постановку, преобладающую в литературе~\cite{cao2018agree, sankar2020groupim}.

\textbf{Гомогенные группы (homogeneous).} Группа формируется так, чтобы попарное сходство всех её членов превышало пороговый квантиль $q_{\mathrm{hom}}$ распределения $\mathrm{sim}(u, v)$ по всем парам пользователей в $\mathcal{U}_{\mathrm{exp}}$:
\begin{equation}
    \forall u, v \in G, \, u \neq v : \mathrm{sim}(u, v) \geq q_{\mathrm{hom}}.
\end{equation}
Процедура: случайный выбор стартового пользователя $u_0$; последовательное добавление пользователей, удовлетворяющих условию, из числа $k$ ближайших соседей по $\mathrm{sim}$ до достижения требуемого размера. Содержательно соответствует ситуации, когда в заведении собрались люди со схожими музыкальными вкусами.

\textbf{Гетерогенные группы (heterogeneous).} Группа формируется так, чтобы попарное сходство было ниже квантиля $q_{\mathrm{het}}$ ($q_{\mathrm{het}} < q_{\mathrm{hom}}$):
\begin{equation}
    \forall u, v \in G, \, u \neq v : \mathrm{sim}(u, v) \leq q_{\mathrm{het}}.
\end{equation}
Процедура аналогична, но добавляются пользователи, наиболее далёкие от уже включённых членов группы (greedy max-min). Это сценарий, в котором гипотеза H3 предсказывает наибольшее преимущество audio-aware агрегатора над классическими стратегиями: при низкой согласованности вкусов простое усреднение предпочтений даёт компромиссные, мало устраивающие всех рекомендации, а адаптивное взвешивание получает наибольшую отдачу.

Чтобы гарантировать различимость трёх сценариев, после генерации проверяется, что эмпирические распределения средних попарных сходств $\overline{\mathrm{sim}}(G) = \binom{|G|}{2}^{-1} \sum_{u, v \in G, u < v} \mathrm{sim}(u, v)$ статистически значимо различны между сценариями (тест Манна--Уитни с поправкой Бонферрони).

\subsection{Размеры групп и протокол оценки}

Размеры групп выбираются из множества $\{2, 4, 8\}$. Эти значения покрывают типичные сценарии совместного потребления: парное прослушивание, малая компания и средняя группа посетителей заведения. Большие размеры ($|G| > 8$) не рассматриваются по двум причинам: (а)~в реальном продуктовом сценарии управление вкусами групп свыше 8~человек выходит за рамки музыкального рекомендера и требует отдельной постановки; (б)~для синтетических групп при больших $|G|$ ограничения на гомогенность/гетерогенность становятся слишком жёсткими и приводят к малому числу допустимых конфигураций, что снижает статистическую мощность оценок.

Каждой группе $G \in \mathcal{G}$ ставится в соответствие \emph{целевой набор треков} $\mathcal{T}_G \subseteq \mathcal{I}$~--- треки, прослушанные хотя бы одним членом группы в тестовом периоде:
\begin{equation}
    \mathcal{T}_G = \bigcup_{u \in G} \left\{ i : (i, t) \in \mathcal{S}_u^{\mathrm{test}} \right\}.
\end{equation}
Альтернативное определение~--- треки, прослушанные \emph{всеми} членами группы~--- для большинства синтетических групп даёт пустое множество и потому не используется. Метрики качества и справедливости (раздел~2.4) вычисляются по $\mathcal{T}_G$ с разделением вкладов отдельных членов: для каждого $u \in G$ метрика рассматривает как релевантные те треки $\mathcal{T}_G$, которые прослушаны именно $u$. Это позволяет определять справедливость как равномерность удовлетворённости членов группы (раздел~2.4) без дополнительных предположений.

Все идентификаторы синтезированных групп, их составы, типы и размеры сохраняются вместе с фиксированными сидами, что обеспечивает воспроизводимость и сопоставимость результатов между разными моделями агрегации.

\section{Метрики оценки}

В настоящем разделе описана система метрик, используемых для проверки сформулированных гипотез. Метрики разделяются на три уровня. \emph{Персональные} метрики применяются на инструментальном этапе при выборе лучшего per-user скорера и проверяют корректность работы воронки candidate generation~$\to$~ranking. \emph{Групповые} метрики оценивают качество предлагаемых агрегаторов и составляют основу проверки центральных гипотез H1--H3; среди них выделяется \emph{главная метрика} DFH@K, объединяющая качество и справедливость. Отдельный блок составляют \emph{метрики cold-start}, оценивающие поведение системы на редких треках.

\subsection{Персональные метрики}

Для каждого пользователя $u \in \mathcal{U}_{\mathrm{exp}}$ модель формирует ранжированный список рекомендаций $\mathrm{Rec}_K(u) = (i_1, i_2, \ldots, i_K)$. Релевантным считается трек, прослушанный пользователем в тестовом периоде: $\mathcal{R}_u = \{ i : (i, t) \in \mathcal{S}_u^{\mathrm{test}} \}$.

\textbf{Recall@K}~--- доля релевантных треков, попавших в верхушку списка:
\begin{equation}
    \mathrm{Recall@K}(u) = \frac{|\mathrm{Rec}_K(u) \cap \mathcal{R}_u|}{|\mathcal{R}_u|}.
\end{equation}

\textbf{HitRate@K}~--- индикатор того, что хотя бы один релевантный трек попал в список:
\begin{equation}
    \mathrm{HR@K}(u) = \mathbf{1}\!\left[ |\mathrm{Rec}_K(u) \cap \mathcal{R}_u| \geq 1 \right].
\end{equation}

\textbf{NDCG@K}~--- нормализованный дисконтированный кумулятивный выигрыш, учитывающий позицию релевантного трека в списке:
\begin{equation}
    \mathrm{DCG@K}(u) = \sum_{k=1}^{K} \frac{\mathbf{1}[i_k \in \mathcal{R}_u]}{\log_2(k + 1)},
    \qquad
    \mathrm{NDCG@K}(u) = \frac{\mathrm{DCG@K}(u)}{\mathrm{IDCG@K}(u)},
\end{equation}
где $\mathrm{IDCG@K}(u)$~--- значение DCG@K при идеальном ранжировании.

\textbf{MRR}~--- средний обратный ранг первого релевантного трека в полном ранжировании каталога:
\begin{equation}
    \mathrm{MRR}(u) = \frac{1}{\min\{k : i_k \in \mathcal{R}_u\}},
\end{equation}
с условием $\mathrm{MRR}(u) = 0$, если $\mathcal{R}_u \cap \mathrm{Rec}(u) = \varnothing$.

Итоговые персональные метрики получаются усреднением по $\mathcal{U}_{\mathrm{exp}}$. Основным значением является $K = 10$; дополнительно репортируются $K = 5$ для строгой оценки верхушки списка и $K = 20$ для согласованности с групповым уровнем (см.~ниже).

\subsection{Групповые метрики}

Для группы $G \in \mathcal{G}$ модель формирует ранжированный список $\mathrm{Rec}_K(G) = (i_1, \ldots, i_K)$. Релевантным для участника $u \in G$ считается трек $i \in \mathcal{R}_u$, как и в персональном случае. Это согласуется с конструкцией целевого набора $\mathcal{T}_G$ из раздела~2.3: групповая рекомендация оценивается через индивидуальные тестовые истории членов группы, а не через искусственно сконструированный <<групповой ground truth>>.

\textbf{Удовлетворённость участника} $\mathrm{Sat}_u(G, K)$ определяется как $\mathrm{NDCG@K}$ персональной выборки $\mathcal{R}_u$ относительно ранжирования $\mathrm{Rec}_K(G)$:
\begin{equation}
    \mathrm{Sat}_u(G, K) = \frac{\sum_{k=1}^{K} \mathbf{1}[i_k \in \mathcal{R}_u] / \log_2(k + 1)}{\mathrm{IDCG@K}(u)}.
\end{equation}
Содержательно $\mathrm{Sat}_u(G, K) \in [0, 1]$ есть мера того, насколько групповой список устраивает конкретного пользователя.

\textbf{Group-NDCG@K} определяется как средняя удовлетворённость по членам группы:
\begin{equation}
    \mathrm{NDCG@K}(G) = \frac{1}{|G|} \sum_{u \in G} \mathrm{Sat}_u(G, K).
\end{equation}

\textbf{Disagreement@K}~--- стандартное отклонение удовлетворённости внутри группы:
\begin{equation}
    \mathrm{Dis@K}(G) = \sqrt{ \frac{1}{|G|} \sum_{u \in G} \left( \mathrm{Sat}_u(G, K) - \mathrm{NDCG@K}(G) \right)^2 }.
\end{equation}
Меньшее значение соответствует более согласованной рекомендации.

\subsection{Основные метрики сравнения: Group-NDCG@K, Jain@K и сводная DFH@K}

В качестве \emph{основных} метрик сравнения моделей агрегации в настоящей работе репортируются \emph{две} величины, отдельно характеризующие качество и справедливость: групповой NDCG@K и индекс Джейна, нормированный в $[0, 1]$ и устойчивый к нулевым значениям~\cite{jin2023fairness_survey}:
\begin{equation}
    \mathrm{Jain}(G, K) = \frac{ \left( \sum_{u \in G} \mathrm{Sat}_u(G, K) \right)^2 }{ |G| \cdot \sum_{u \in G} \mathrm{Sat}_u(G, K)^2 }.
\end{equation}
Значение $\mathrm{Jain} = 1$ соответствует абсолютно равномерной удовлетворённости, $\mathrm{Jain} = 1/|G|$~--- ситуации, в которой удовлетворён ровно один участник.

Дополнительно вводится \emph{сводная} метрика \emph{Dual Fairness-Harmonic at K}, определяемая как гармоническое среднее группового NDCG@K и индекса Джейна:
\begin{equation}
    \mathrm{DFH@K}(G) = \frac{2 \cdot \mathrm{NDCG@K}(G) \cdot \mathrm{Jain}(G, K)}{ \mathrm{NDCG@K}(G) + \mathrm{Jain}(G, K) }.
\end{equation}
DFH@K служит компактной сводной мерой и используется как критерий выбора лучшей конфигурации гиперпараметров на валидации. Однако приоритет в интерпретации результатов отдаётся раздельному анализу Group-NDCG@K и Jain@K по двум методологическим причинам. Во-первых, типичные значения Group-NDCG@K в задачах рекомендаций лежат в диапазоне $0{,}05$--$0{,}25$, тогда как Jain@K~--- в диапазоне $0{,}6$--$1{,}0$; гармоническое среднее величин разного масштаба смещено в сторону меньшей компоненты, что ослабляет вклад Jain в DFH в данном диапазоне. Во-вторых, раздельный отчёт качества и справедливости содержательно богаче и позволяет интерпретировать источник различий между агрегаторами~--- выигрыш по качеству, по справедливости либо по обеим компонентам одновременно.

Итоговые значения метрик по сценарию агрегируются по группам в пределах каждой комбинации $(\mathrm{type}(G), |G|)$ и сопровождаются бутстрап-доверительными интервалами уровня $0{,}95$ ($N_{\mathrm{boot}} = 1000$).

\subsection{Метрики cold-start}

Для содержательного анализа того, насколько преимущество audio-aware агрегатора связано с его способностью работать на редких треках (вспомогательный вопрос к гипотезе~H1), вводится отдельная подвыборка треков с низкой частотой взаимодействий. Трек $i$ относится к cold-start множеству $\mathcal{I}_{\mathrm{cold}}$, если число его прослушиваний в обучающем периоде не превышает фиксированный порог $n_{\mathrm{cold}}$:
\begin{equation}
    \mathcal{I}_{\mathrm{cold}} = \left\{ i \in \mathcal{I} : \sum_{u \in \mathcal{U}_{\mathrm{exp}}} \mathbf{1}\!\left[ i \in \mathcal{S}_u^{\mathrm{train}} \right] \leq n_{\mathrm{cold}} \right\}.
\end{equation}
Использование фиксированного порога предпочтительнее квантильного: оно делает условие воспроизводимым между разными подвыборками и упрощает сопоставление результатов. Конкретное значение $n_{\mathrm{cold}}$ выбирается так, чтобы $|\mathcal{I}_{\mathrm{cold}}|$ составляло заметную долю каталога ($\sim 10$--$20\%$) и обеспечивало статистически устойчивые оценки.

На cold-start подвыборке вычисляются те же персональные и групповые метрики, что и в основном протоколе, но релевантные множества переопределяются как $\mathcal{R}_u^{\mathrm{cold}} = \mathcal{R}_u \cap \mathcal{I}_{\mathrm{cold}}$, а ранжирование ограничивается треками из $\mathcal{I}_{\mathrm{cold}}$. Это даёт прямую количественную оценку способности моделей рекомендовать редкие треки, не доминируемые популярными элементами каталога.

\subsection{Сводка отчётных метрик}

В итоговой таблице результатов для каждой комбинации (модель агрегации, $\mathrm{type}(G)$, $|G|$) репортируются основные метрики Group-NDCG@K и Jain@K при $K = 10$ и $K = 20$, сводная DFH@K, вспомогательная Disagreement@K, а также Group-NDCG@K и Jain@K на cold-start подвыборке. Персональные метрики per-user скорера (Recall@K, NDCG@K, HR@K, MRR) репортируются справочно в приложении к работе для контроля качества зафиксированного $f_\theta^\ast$. Все метрики сопровождаются доверительными интервалами уровня $0{,}95$, оценёнными бутстрапом по группам ($N_{\mathrm{boot}} = 1000$ итераций).

\section{Экспериментальный протокол: персональный уровень}

Настоящий раздел описывает протокол подготовки зафиксированного per-user скорера $f_\theta^\ast$, выступающего неизменным входом группового агрегатора. Воронка состоит из трёх инженерных этапов; ablation между альтернативными архитектурами на персональном уровне в основном тексте работы \emph{не проводится}, поскольку выбор скорера не является предметом исследования. Цель раздела~--- зафиксировать процедуру получения качественных персональных скоров и кандидатных множеств, на основе которых строится групповой эксперимент.

\subsection{Этап 1: генерация кандидатов ($|\mathcal{I}| \to 10^4$)}

Модель этапа~--- ALS (Alternating Least Squares)~\cite{rendle2009bpr}: матричная факторизация на матрице неявных предпочтений, представляющая каждого пользователя и трек векторами фиксированной размерности $d_{\mathrm{ALS}}$. Скор вычисляется как $\hat{s}_{u, i}^{(1)} = \mathbf{p}_u^\top \mathbf{q}_i$. Для эффективного отбора $K_1 = 10^4$ кандидатов на пользователя строится ANN-индекс HNSW~\cite{johnson2019faiss} по векторам треков $\{\mathbf{q}_i\}_{i \in \mathcal{I}}$. К выходу этапа подмешиваются два safety-net набора: топ-$N_{\mathrm{pop}}$ по абсолютной популярности и топ-$N_{\mathrm{rec}}$ по recency-скорам в пределах последнего временного окна. Объединённое множество кандидатов обозначается $\mathcal{C}_u^{(1)}$.

Этап~1 не подвергается ablation. Корректность работы проверяется единственным инженерным условием: $\mathrm{Recall@}K_1 \geq 0{,}95$ на валидационной выборке. При невыполнении условия пробуется альтернатива (BPR-MF) или увеличение $K_1$.

\subsection{Этап 2: лёгкое pre-ranking ($10^4 \to 10^3$)}

Модель этапа~--- лёгкий CatBoost-классификатор~\cite{prokhorenkova2018catboost} на узком наборе дешёвых признаков:
\begin{itemize}
    \item скор ALS $\hat{s}_{u, i}^{(1)}$;
    \item статистики популярности трека (общая, recency-взвешенная);
    \item статистики истории пользователя ($|\mathcal{S}_u|$, доля повторных прослушиваний);
    \item косинусное сходство аудиопрофиля пользователя $\bar{\mathbf{a}}_u$ с аудиоэмбеддингом трека $\mathbf{a}_i$ (один скаляр, не плотный вектор).
\end{itemize}
Полные плотные аудиоэмбеддинги $\mathbf{a}_i$ на этом этапе не используются: при $|\mathcal{C}_u^{(1)}| \approx 10^4$ это привело бы к существенному росту стоимости обучения без значимого выигрыша в качестве отбора кандидатов. Этап~2 также не подвергается ablation; инженерный критерий приёмки~--- $\mathrm{Recall@}K_2 \geq 0{,}92$ на валидации, $K_2 = 10^3$.

\subsection{Этап 3: последовательное ранжирование ($10^3 \to 10^2$)}

Модель этапа зафиксирована как gSASRec~\cite{gsasrec2023}~--- трансформерная архитектура с каузальной маской и функцией потерь gBCE, компенсирующей смещение от отрицательного семплирования. Выбор обоснован выше (раздел~2.1.1): согласно литературе, gSASRec систематически превосходит SASRec~\cite{sasrec2018} и BERT4Rec~\cite{bert4rec2019} на больших каталогах, к которым относится YAMBDA. Сравнение с альтернативными архитектурами (Mamba4Rec~\cite{liu2024mamba4rec}) представляет самостоятельный интерес и при необходимости выносится в приложение работы как инструментальный этап подбора скорера.

Модель обучается на множестве кандидатов $\mathcal{C}_u^{(2)}$ этапа~2, а не на полном каталоге, что согласуется с её ролью ранжировщика, а не генератора кандидатов. Целевой элемент~--- следующий трек в обучающей последовательности; для каждого позитивного примера семплируется $n_{\mathrm{neg}}$ отрицательных из $\mathcal{C}_u^{(2)} \setminus \mathcal{S}_u^{\mathrm{train}}$.

Выходом этапа~3 является топ-$K_3 = 10^2$ кандидатов на пользователя с соответствующими скорами $\hat{s}_{u, i}^{(3)}$. Эти кандидаты и скоры составляют зафиксированный $f_\theta^\ast$, на котором строится групповой эксперимент. Кэширование выходов ($\mathcal{C}_u^{(3)}$ и $\hat{s}_{u, i}^{(3)}$) выполняется однократно для всех $u \in \mathcal{U}_{\mathrm{exp}}$, что исключает повторные прогоны нейросетевой модели при каждом запуске группового эксперимента (раздел~2.7).

\subsection{Подбор гиперпараметров и ранняя остановка}

Подбираемые гиперпараметры gSASRec включают размерность эмбеддингов $d$, максимальную длину последовательности $L_{\max}$, число слоёв и голов внимания, параметр калибровки $\alpha$ функции gBCE, начальный learning rate и dropout. Для CatBoost-этапа~2 подбираются глубина деревьев, learning rate, число итераций и $\ell_2$-регуляризация. Размерность ALS $d_{\mathrm{ALS}}$ выбирается из соображений приёмки этапа~1.

Поиск выполняется случайным сэмплированием с фиксированным бюджетом обучения. Лучшая конфигурация выбирается по NDCG@10 на валидационной выборке. Ранняя остановка применяется по тому же критерию с терпением $p_{\mathrm{patience}} = 5$ эпох. Все случайные операции используют фиксированные сиды, включая отдельную фиксацию состояния CUDA-генератора. Итоговые метрики на тесте сообщаются как среднее по $n_{\mathrm{seeds}} = 5$ запускам с разными сидами и приводятся справочно в приложении.

Зафиксированный $f_\theta^\ast$ далее является неизменным входом всех экспериментов с групповыми агрегаторами. Это разделение исключает смешение источников вариации: различия между агрегаторами не могут быть объяснены различиями в качестве per-user скорера.

\section{Экспериментальный протокол: групповой уровень}

Настоящий раздел описывает протокол обучения и сравнения групповых агрегаторов~--- центральный экспериментальный блок работы. На вход всех агрегаторов подаются персональные скоры зафиксированного $f_\theta^\ast$, обучаемые компоненты~--- только параметры самого агрегатора. Такое разделение обеспечивает, что наблюдаемые различия в групповых метриках объясняются исключительно стратегией агрегации.

\paragraph{Точка подключения в воронке.} Агрегатор подключается \emph{непосредственно после} этапа~3 персональной воронки (раздел~2.5). Для каждого члена группы $u \in G$ независимо выполняются этапы~1--3, дающие топ-$K_3 = 10^2$ кандидатов с персональными скорами $\hat{s}_{u, i}^{(3)}$. Агрегатор объединяет персональные кандидатные множества членов группы в общий пул $\mathcal{C}_G = \bigcup_{u \in G} \mathcal{C}_u^{(3)}$ и приписывает каждому треку $i \in \mathcal{C}_G$ групповой скор $\hat{s}_{G, i}$. Итоговый список $\mathrm{Rec}_K(G)$ формируется как $\mathrm{argtop}\text{-}K$ по групповому скору на пуле $\mathcal{C}_G$. Финальный CatBoost-этап (этап~4 архитектуры из главы~1) в техническом эксперименте \emph{не применяется}: его обучение проводится на распределении персональных скоров и индивидуальных таргетах, тогда как на инференсе он получает агрегированный групповой скор на агрегированном пуле кандидатов. Это создаёт дополнительный, не относящийся к предмету исследования источник вариации между агрегаторами; в продуктовой части диплома этап~4 переобучается на групповом распределении.

\subsection{Состав сравниваемых агрегаторов}

Сравниваются восемь стратегий, разбитых на три блока.

\textbf{Классические необучаемые стратегии~\cite{cao2018agree}.}
\begin{itemize}
    \item \textit{Average} (AVG): $\hat{s}_{G, i} = \frac{1}{|G|} \sum_{u \in G} \hat{s}_{u, i}$.
    \item \textit{Least Misery} (LM): $\hat{s}_{G, i} = \min_{u \in G} \hat{s}_{u, i}$.
    \item \textit{Max Pleasure} (MP): $\hat{s}_{G, i} = \max_{u \in G} \hat{s}_{u, i}$.
\end{itemize}
Эти стратегии не имеют обучаемых параметров и применяются непосредственно поверх $f_\theta^\ast$.

\textbf{Обучаемые бейзлайны для эфемерных групп.}
\begin{itemize}
    \item \textit{ID-based AGREE}~\cite{cao2018agree}~--- attention-агрегатор, использующий обучаемые ID-эмбеддинги треков $\mathbf{e}_i \in \mathbb{R}^{d_e}$:
    \begin{equation}
        \alpha_{u, i}^{(G)} = \mathrm{softmax}_{u \in G}\, \phi_\psi(\mathbf{e}_u, \mathbf{e}_i),
    \end{equation}
    где $\phi_\psi$~--- двухслойный перцептрон с активацией ReLU. Включается как стандартная attention-точка сравнения.
    \item \textit{GroupIM}~\cite{sankar2020groupim}~--- модель максимизации взаимной информации между представлением группы и индивидуальными представлениями членов, специально разработанная для эфемерных групп. Её включение принципиально: GroupIM~--- наиболее релевантный современный бейзлайн для постановки настоящей работы (синтетические эфемерные группы), что делает сравнение с ним обязательным для содержательной валидации гипотезы H1. Используется реализация авторов с адаптацией входа под пул $\mathcal{C}_G$ и зафиксированный $f_\theta^\ast$ (детали~--- в приложении).
\end{itemize}

\textbf{Предложенное семейство audio-aware агрегаторов.} Три варианта, описанные в разделе~2.1: \textit{Audio taste profile} (необучаемый), \textit{Audio-AGREE} (обучаемый attention поверх $\bar{\mathbf{a}}_u$ и $\mathbf{a}_i$), \textit{group cross-attention} (multi-head cross-attention с членами группы как ключами и значениями). В audio-aware вариантах ID-эмбеддинги треков не используются~--- агрегатор оперирует исключительно аудиоэмбеддингами $\mathbf{a}_i$ из YAMBDA.

\subsection{Обучающие группы и протокол семплирования}

Обучаемые агрегаторы (ID-based AGREE, GroupIM, Audio-AGREE, group cross-attention) тренируются на синтетических группах, сформированных исключительно из обучающего периода. Целевые группы оценочного протокола (раздел~2.3) при обучении не используются~--- это исключает любые формы утечки тестовых данных.

Принципиальное методологическое решение~--- использование при обучении только \emph{случайных} групп, без сценариев гомогенности и гетерогенности. Обоснование двоякое. С методологической стороны: совпадение распределений обучающих и оценочных групп по характеристикам сходства привело бы к ситуации, в которой выигрыш модели объясняется подгонкой под распределение, а не качеством архитектуры. С продуктовой стороны: в реальном сценарии заведения структура группы заранее неизвестна, и модель должна обобщать на произвольные конфигурации сходства членов.

Группы семплируются \emph{онлайн} на каждой эпохе обучения: для каждого батча генерируются новые группы случайных пользователей с размерами $|G| \in \{2, 4, 8\}$, распределёнными равномерно. Объём обучающих групп за эпоху составляет $N_{\mathrm{train}}^G$ (типичное значение~--- $5 \cdot 10^4$); общее число эпох регулируется ранней остановкой. Онлайн-семплирование обеспечивает экспозицию модели к большему разнообразию групповых конфигураций без раздувания памяти.

Для каждой обучающей группы $G$ целевое множество определяется как объединение следующих треков членов группы в обучающем периоде:
\begin{equation}
    \mathcal{Y}_G = \bigcup_{u \in G} \mathcal{R}_u^{\mathrm{next}}, \qquad \mathcal{R}_u^{\mathrm{next}} = \left\{ i : (i, t) \in \mathcal{S}_u^{\mathrm{train}}, \, t \in W_u \right\},
\end{equation}
где $W_u$~--- скользящее окно следующих $n_W$ событий пользователя $u$, не используемых при формировании контекста. Использование объединения, а не пересечения, согласовано с конструкцией оценочного множества $\mathcal{T}_G$ (раздел~2.3) и обусловлено тем, что для большинства синтетических групп пересечение пусто.

\subsection{Функция потерь и режимы обучения}

Основной режим обучения~--- pairwise BPR-loss на объединении. Для каждой пары $(G, i^+)$, где $i^+ \in \mathcal{Y}_G$, семплируется отрицательный пример $i^- \notin \mathcal{Y}_G$ из распределения, пропорционального популярности (popularity-based negative sampling), и оптимизируется:
\begin{equation}
    \mathcal{L}_{\mathrm{rank}} = - \sum_{(G, i^+, i^-)} \log \sigma\!\left( \hat{s}_{G, i^+} - \hat{s}_{G, i^-} \right).
\end{equation}
Pairwise BPR выбран как компромисс между точечной BCE (плохо учитывает порядок) и listwise softmax (требует дорогостоящего семплирования при $|\mathcal{I}| \approx 9{,}4 \cdot 10^6$). На каждый позитивный пример используется $n_{\mathrm{neg}}^G$ отрицательных, что эффективно даёт листовую оценку при ограниченной стоимости.

Дополнительно вводится \emph{fairness-aware режим} как отдельная конфигурация эксперимента:
\begin{equation}
    \mathcal{L}_{\mathrm{fair}} = \mathcal{L}_{\mathrm{rank}} + \lambda \cdot \mathrm{Var}_{u \in G}\!\left( \hat{s}_{u, i^+} \cdot \alpha_{u, i^+}^{(G)} \right),
\end{equation}
где штрафной член компенсирует неравномерность вклада членов группы в положительные предсказания. Использование fairness-aware потерь методологически нетривиально: оптимизация модели по компоненте, явно входящей в оценочную метрику DFH@K, повышает риск переоценить её преимущество, и это явно оговаривается при представлении результатов. Соответственно, fairness-aware режим не используется для сравнения с классическими бейзлайнами по DFH@K~--- его роль в работе ограничена анализом того, насколько прирост справедливости достигается естественным образом из rank-loss и насколько~--- за счёт явной регуляризации.

Per-user скорер $f_\theta^\ast$ во всех режимах остаётся полностью замороженным: обучаются исключительно параметры $\psi$ агрегатора. Это~--- единственный способ обеспечить, чтобы сравнение разных агрегаторов отражало именно стратегию агрегации, а не косвенное дообучение нижележащего скорера.

\subsection{Гиперпараметры и ранняя остановка}

Подбираемые гиперпараметры обучаемых агрегаторов:
\begin{itemize}
    \item размерности скрытых слоёв $\phi_\psi$ и количество голов внимания (для group cross-attention);
    \item dropout, начальный learning rate, размер батча групп $B^G$;
    \item число отрицательных семплов на позитивный $n_{\mathrm{neg}}^G$;
    \item для fairness-aware режима~--- коэффициент $\lambda$;
    \item для Audio taste profile~--- единственный параметр, температура $\tau$.
\end{itemize}
Поиск выполняется случайным сэмплированием с фиксированным бюджетом обучения, единым для всех обучаемых агрегаторов. Лучшая конфигурация выбирается по DFH@10 на валидационных группах~--- сгенерированных из обучающего периода по той же процедуре, что и тестовые, но с отдельным сидом и из непересекающегося подмножества пользователей и временного окна.

Ранняя остановка применяется по DFH@10 на валидационных группах с терпением $p_{\mathrm{patience}}^G = 3$ эпохи; верхний предел эпох~--- $E_{\max}^G$. Итоговые результаты на тестовых группах сообщаются как среднее по $n_{\mathrm{seeds}}^G = 5$ запускам с разными случайными сидами.

\subsection{Протокол оценки}

После обучения каждый агрегатор оценивается на тестовых группах $\mathcal{G}$ (раздел~2.3). Для каждой группы $G$ вычисляется ранжирование $\mathrm{Rec}_K(G)$ путём прогона зафиксированного $f_\theta^\ast$ для всех её членов и применения агрегатора к получаемым скорам. Метрики из раздела~2.4 (DFH@K, Group-NDCG@K, Jain@K, Disagreement@K) усредняются по группам в пределах каждой комбинации $(\mathrm{type}(G), |G|)$ и репортируются с бутстрап-доверительными интервалами уровня $0{,}95$.

Сравнение моделей проводится в трёх срезах:
\begin{itemize}
    \item \textbf{общий срез} (все группы независимо от типа и размера) для проверки гипотезы H1: audio-aware семейство против AVG, LM, MP, ID-based AGREE и GroupIM;
    \item \textbf{срез по архитектурной сложности} (audio taste profile vs Audio-AGREE vs group cross-attention) для проверки H2;
    \item \textbf{дополнительные содержательные срезы} (тип группы, cold-start подвыборка) для интерпретации источников выигрыша audio-aware агрегатора~--- без оформления самостоятельной гипотезы.
\end{itemize}
Статистическая значимость попарных различий между моделями устанавливается парным $t$-критерием по группам с поправкой Бонферрони на множественность сравнений в рамках каждого среза.

\section{Вычислительные ресурсы и воспроизводимость}

\subsection{Среда выполнения}

Все эксперименты выполняются в среде Google Colab с подключением графических ускорителей NVIDIA A100 (40~ГБ памяти) на ограниченное по продолжительности время. Доступная оперативная память хоста~--- 80~ГБ. Эти ограничения определяют ряд практических решений, описанных ниже.

В качестве программного стека используются: \texttt{Python 3.11}, \texttt{PyTorch} в качестве основной фреймворковой библиотеки для обучения нейросетевых моделей, \texttt{CatBoost} для ранжирующей модели, \texttt{NumPy}, \texttt{pandas} и \texttt{Polars} для предварительной обработки данных, \texttt{scikit-learn} для кластеризации и расчёта вспомогательных статистик, \texttt{SciPy} для статистических тестов. Версии всех зависимостей фиксируются в файле \texttt{requirements.txt}, сохраняемом в репозитории работы.

\subsection{Управление вычислительной стоимостью}

Полный объём YAMBDA ($4{,}79 \cdot 10^9$ событий) превышает доступные вычислительные ресурсы на несколько порядков. Применяемые ниже меры позволяют свести стоимость обучения и оценки к допустимым значениям без потери методологической корректности.

\textbf{Подвыборка пользователей и треков.} Стратифицированная подвыборка $\mathcal{U}_{\mathrm{exp}}$ из раздела~2.2.4 ограничивает объём событий до уровня, при котором обучение каждой модели-кандидата занимает не более нескольких часов на A100.

\textbf{Кэширование промежуточных представлений.} Топ-$K_3 = 100$ кандидатов и соответствующие скоры $\hat{s}_{u, i}^{(3)}$ зафиксированной модели $f_\theta^\ast$ (выход этапа~3 воронки) вычисляются однократно для всех $u \in \mathcal{U}_{\mathrm{exp}}$ и сохраняются на диск. На этапе обучения и оценки групповых агрегаторов используется этот кэш, что исключает повторные прогоны нейросетевой модели на каждом батче и многократно сокращает время эксперимента.

\textbf{Кэширование промежуточных этапов.} Аналогично кэшируются выходы этапа~1 (топ-$K_1$ кандидатов ALS) и этапа~2 (топ-$K_2$ после лёгкого CatBoost), что исключает повторные пересчёты при подборе гиперпараметров gSASRec на этапе~3.

\textbf{Смешанная точность (mixed precision).} Обучение нейросетевых моделей выполняется в \texttt{bfloat16} на матричных операциях с сохранением \texttt{float32} в накопителях градиентов, что снижает потребление памяти и ускоряет обучение без заметного влияния на качество.

\subsection{Воспроизводимость}

Воспроизводимость экспериментов обеспечивается на нескольких уровнях.

\textbf{Фиксация случайных операций.} Все случайные генераторы~--- \texttt{Python random}, \texttt{NumPy}, \texttt{PyTorch} (CPU и CUDA)~--- инициализируются явно заданными сидами. Для операций CUDA, чувствительных к недетерминизму (свёртки, atomic operations), активируются флаги детерминированного выполнения; в случаях, где это невозможно по производительности, недетерминизм компенсируется усреднением по $n_{\mathrm{seeds}}$ запускам.

\textbf{Версионирование данных.} Идентификаторы пользователей и треков, вошедших в $\mathcal{U}_{\mathrm{exp}}$, составы синтезированных групп, точки временного сплита $t_{\mathrm{train}}, t_{\mathrm{val}}$ и список cold-start треков $\mathcal{I}_{\mathrm{cold}}$ сохраняются в виде CSV/Parquet файлов, версионируемых вместе с кодом. Это позволяет любому третьему лицу воспроизвести экспериментальные выборки от исходного датасета YAMBDA до итоговых метрик.

\textbf{Управление конфигурациями.} Все гиперпараметры экспериментов задаются в декларативных конфигурационных файлах формата YAML, привязанных к git-коммитам. Каждый запуск эксперимента порождает запись в журнале с указанием хеша коммита, конфигурации, использованных сидов и результирующих метрик.

\textbf{Публикация кода.} Исходный код реализации, скрипты обработки данных, конфигурации экспериментов и журналы запусков публикуются в открытом репозитории с лицензией, допускающей независимое воспроизведение и расширение.

\subsection{Ограничения и план их компенсации}

Совокупность вычислительных ограничений накладывает на работу несколько методологических ограничений, которые открыто фиксируются.

Во-первых, эксперименты ведутся на подвыборке пользователей $\mathcal{U}_{\mathrm{exp}}$, а не на полном объёме YAMBDA. Это компенсируется стратификацией подвыборки и оценкой устойчивости результатов к выбору сида подвыборки на репрезентативной части экспериментов.

Во-вторых, число эпох обучения и количество запусков с разными сидами ($n_{\mathrm{seeds}} = 5$) ограничены доступным временем и являются нижней границей для устойчивых статистических выводов. Все статистические тесты явно учитывают это число запусков.

В-третьих, гиперпараметрический поиск выполняется случайным сэмплированием с фиксированным бюджетом, а не полным перебором или байесовской оптимизацией. Это создаёт риск неоптимальной конфигурации для отдельных моделей; риск контролируется единым бюджетом для всех моделей и явным репортом всего пройденного пространства гиперпараметров в приложении к работе.

\section{Выводы по главе}

В настоящей главе сформулирован экспериментальный план, сознательно сужённый по сравнению с архитектурным ориентиром главы~1 ради методологической чистоты сравнения групповых агрегаторов.

\textbf{Постановка задачи и гипотезы.} Центром экспериментального исследования является уровень групповой агрегации. Предложено семейство audio-aware attention-агрегаторов, использующих аудиоэмбеддинги YAMBDA вместо обучаемых ID-эмбеддингов. Сформулированы две центральные гипотезы: H1~--- преимущество предложенного семейства над классическими стратегиями (AVG, LM, MP) и современными обучаемыми бейзлайнами для эфемерных групп (ID-based AGREE, GroupIM); H2~--- монотонный рост качества по архитектурной сложности audio-aware агрегатора. Дополнительные срезы (тип группы, cold-start) обеспечивают содержательную интерпретацию источников выигрыша, но не оформляются как самостоятельные гипотезы. Поведенческая валидация аудиоэмбеддингов проводится как методологическая sanity-проверка перед основным экспериментом.

\textbf{Данные.} Описана подготовка YAMBDA: фильтрация по типу взаимодействий и флагу \texttt{is\_organic}, пороги активности, сохранение повторных прослушиваний, глобальный временной сплит вместо leave-one-out, стратифицированная подвыборка под вычислительные ограничения. Отдельным подразделом введена поведенческая валидация аудиоэмбеддингов как замена недоступной семантической верификации.

\textbf{Синтез групп.} Из индивидуальных профилей формируются эфемерные группы трёх типов (random, homogeneous, heterogeneous) трёх размеров ($|G| \in \{2, 4, 8\}$). Сходство пользователей определяется через коэффициент Жаккара по множествам прослушанных треков~--- сознательно не через аудиоэмбеддинги, чтобы исключить циркулярность с предлагаемым агрегатором.

\textbf{Метрики.} Основными метриками сравнения выступают Group-NDCG@K и Jain@K, репортируемые раздельно; сводная DFH@K используется как вспомогательный композитный показатель и критерий выбора гиперпараметров на валидации. Cold-start подвыборка треков формируется по фиксированному порогу частоты и оценивается отдельно. Все метрики сопровождаются бутстрап-доверительными интервалами уровня $0{,}95$.

\textbf{Протокол персонального уровня.} Реализуется трёхэтапная воронка: ALS с HNSW-индексом ($|\mathcal{I}| \to 10^4$), лёгкий CatBoost ($10^4 \to 10^3$), последовательное ранжирование gSASRec ($10^3 \to 10^2$). Финальный этап тяжёлого CatBoost (этап~4 архитектуры главы~1) в групповом эксперименте \emph{не применяется}, чтобы избежать сдвига распределений между его обучением и инференсом; его реализация относится к продуктовой части диплома. Архитектура этапа~3 зафиксирована как gSASRec без ablation против альтернатив. Выход этапа~3~--- зафиксированный $f_\theta^\ast$.

\textbf{Протокол группового уровня.} Сравниваются восемь стратегий: AVG, LM, MP, ID-based AGREE, GroupIM и три варианта audio-aware агрегатора (taste profile, Audio-AGREE, group cross-attention). Обучаемые модели тренируются на синтетических случайных группах из обучающего периода с pairwise BPR-функцией потерь и popularity-based негативным семплированием; целевое множество~--- объединение следующих треков членов группы. Per-user скорер $f_\theta^\ast$ во всех конфигурациях заморожен. Итоговый список $\mathrm{Rec}_K(G)$ формируется как $\mathrm{argtop}\text{-}K$ по групповому скору на пуле $\mathcal{C}_G$, без последующего CatBoost-этапа. Сравнение проводится в трёх срезах с парным $t$-критерием и поправкой Бонферрони.

\textbf{Воспроизводимость.} Все случайные операции, выборки данных и конфигурации фиксируются и версионируются. Кэширование скоров $f_\theta^\ast$, усечение кандидатов и смешанная точность позволяют выполнить полный экспериментальный план в пределах вычислительных ограничений Google Colab.

Описанный план обеспечивает однозначную проверку гипотез H1 и H2 при минимальном числе движущихся частей и формирует основу для главы~3, в которой представлены результаты экспериментов и их интерпретация.
